%==============================================================================
\section{Workload Model}
\label{sec:model}
%==============================================================================

We study a generic \emph{batched update-query workload} parameterized by three integers:

\begin{itemize}
  \item $n$: the number of items in the collection.
  \item $k$: the number of items whose sort keys change per \emph{epoch}.
    Positions of updated items are known to the caller.
  \item $Q$: the number of \emph{indexed rank queries} per epoch.
    A rank query $\texttt{get\_rank}(i)$ retrieves the item currently at
    sorted position $i$.
\end{itemize}

An \emph{epoch} proceeds as follows: (1) $k$ items are updated at known indices;
(2) the sorted order is restored; (3) $Q$ rank queries are answered.
Steps (1) and (2) alternate with step (3) repeatedly.

\paragraph{Data structures.}
We compare two canonical choices:

\begin{description}
  \item[Sorted Vec + \DeltaSort.]
    Maintains a contiguous sorted array.
    Step (2) costs $T_{\Delta}(n,k) = O(n\sqrt{k})$ in expectation.
    Step (3) costs $O(1)$ per query (direct index).
    Total per epoch: $T_{\Delta}(n,k) + O(Q)$.

  \item[BTreeMap.]
    A self-balancing search tree (Rust's \texttt{std::collections::BTreeMap}).
    Step (2): each updated item is removed and reinserted in $O(\log n)$;
    total $O(k \log n)$.
    Step (3): the item at rank $i$ requires an in-order walk of $i$ nodes,
    costing $O(i)$.
    For $Q$ rank queries with uniformly random ranks, expected cost is
    $Q \cdot O(n/2) = O(Qn)$.
    Alternatively, the full sorted order can be materialized once in $O(n)$
    and all $Q$ queries answered in $O(1)$ each; choose the cheaper option.
    Total per epoch: $O(k \log n + \min(Qn,\; n + Q))$.
\end{description}

This model captures, for example, a leaderboard (updates from score changes,
queries for top-$K$ ranking), a MIDI event list (updates from note drags,
queries for playback iteration), or a sorted search result cache (updates
from index refreshes, queries from UI rendering).
